\scanpage{56}
% archive sheet number: 55

\handwritten{Suppose. first $k$ algebraically closed, $G$ reduced and connected and $H$ reduced. According to (\uline{BIBLE} \deleted{\unsure{VII|VIII}}\inserted{\deleted{Lecture} 4, corollary of theorem \uline{1}\begin{tikzpicture}[x=1cm,y=1cm,baseline=-.1cm]\draw[thin] (0,-.1).. controls (-.1,-.35) and (-1.5,-.1)..(-1.7,-.4)--(-1.75,-.6)--(-1.8,-.35);\end{tikzpicture}}), there exists a finite-dimensional linear representation $\rho:G\longrightarrow\mathrm{GL}(V)$ such that $(\operatorname{Ker}\rho)_{\mathrm{red}}=H$. By 11.11, $G/\operatorname{Ker}\rho$ is affine. hence $(G/H)/(\operatorname{Ker}\rho/H)$ is affine, hence also $G/H$, the group $\operatorname{Ker}\rho/H$ being finite.}

\handwritten{If $k$ is algebraically closed, and $G$ and $H$ reduced, then $(G/H)/(G^{\unsure{\circ|0|o}}/H\cap G^{\unsure{\circ|0|o}})$ is finite,\inserted{ as a quotient of $G/G^{\unsure{\circ|0|o}}$, and\begin{tikzpicture}[x=1cm,y=1cm,baseline=-.1cm]\draw[thin] (0,.05).. controls (0,-.25) and (-1.9,-.15)..(-2,-.35)--(-2.02,-.55);\end{tikzpicture}} $G^{\unsure{\circ|0|o}}/H\cap G^{\unsure{\circ|0|o}}$ is affine, hence $G/H$ is affine.}
\margin{\figure{\handwritten{\begin{tikzpicture}[x=0.2cm,y=0.4cm]\draw (0,8) .. controls (1.2,6.8) and (1,3) .. (1.2,0);\draw (0.8,6.8).. controls (0.7,6.5) and (0.3,6.5)..(-0.2,6.45);\end{tikzpicture}}}}

\figure{\handwritten{\begin{tikzpicture}\draw (0,0) .. controls (3,0.2) and (8,0.2) .. (12,0);\end{tikzpicture}}}
\deleted{\handwritten{\uline{N.B.}1) the theorem \unsure{implies|implies} that if $G$ is affine algebraic and $H$ normal, there exists a finite-dimensional lin\unsure{e|e}ar representation $\rho:G\longrightarrow\mathrm{GL}(V)$ such that $H=\operatorname{Ker}\rho$}}

\deleted{\handwritten{2) The reference to BIBLE is a corollary of 11.16, the d\unsure{e|e}monstration is easy but a little long.}}

\deleted{\handwritten{3) In the reference to BIBLE Suppose $G$ connected, for the demonstration in BIBLE is \unsure{formulated} in the general case\begin{tikzpicture}[x=1cm,y=1cm,baseline=-.1cm]\draw[thin] (0,0).. controls (.45,.12) and (.85,.6)..(1.22,.5).. controls (1.43,.47) and (1.48,.25)..(1.4,.16);\end{tikzpicture}}}
\figure{\handwritten{\begin{tikzpicture}\draw (0,0) .. controls (4,0.4) and (8,0.5) .. (12,0.6);\end{tikzpicture}}}
